\documentclass{article}
\usepackage[T1]{fontenc}
\usepackage{lmodern}
\usepackage{graphicx}
\usepackage{amsmath,amssymb}
\usepackage[a4paper,margin=2cm]{geometry}
\usepackage[absolute,overlay]{textpos}
\setlength{\TPHorizModule}{1mm}
\setlength{\TPVertModule}{1mm}
\usepackage[indonesian]{babel}
\usepackage{hyperref}
\setlength{\emergencystretch}{2em}
% unit: Halaman 56

\begin{document}

% === Halaman 56 (python/native) ===

2. Diffusion

• Prinsip ini menyebarkan pengaruh satu bit plainteks atau kunci ke sebanyak mungkin bit-

bit cipherteks. 

    Contoh: satu huruf plainteks mempengaruhi nilai banyak huruf cipherteks


• Sebagai efek difusi, maka pengubahan kecil pada plainteks sebanyak satu atau dua bit 

menghasilkan perubahan pada bit-bit cipherteks yang tidak dapat diprediksi. 

• Contoh: mengenkripsi plainteks 1111111111111111 

    menghasilkan cipherteks 0110110000101001. 

    Misalkan 1 bit pertama plainteks diubah menjadi 0111111111111111. Jika difusi bekerja, 

maka cipherteksnya berubah signifikan menjadi  1011001011101111

• Diffusion did alam block cipher dapat direalisasikan dengan menggunakan teknik 

permutasi atau transposisi secara berulang-ulang. Permutasi adalah mengacak susunan 

bit atau byte di dalam plainteks.

56

\end{document}
